% !PRINTED 435
\weberpage{435}
%<<CONT
\weq{10}{\begin{gathered} x A_n(x) - A_{n-1}(x) \\ = \overset{\nu}{\Sigma}(-1)^\nu x^{n-2\nu+1} \dfrac{n B_\nu^{(n-\nu)} + (n-1) B_{\nu-1}^{(n-\nu)}}{n-\nu}, \quad 0 \leqq \nu \leqq \dfrac{n+1}{2}. \end{gathered}}
Nun ist aber, wie sich aus den Ausdrücken für die Binomialco\"efficienten leicht ergiebt (\S.\,8),
\[
\dfrac{n B_\nu^{(n-\nu)} + (n-1) B_{\nu-1}^{(n-\nu)}}{n-\nu} = \dfrac{n+1}{n-\nu+1} B_\nu^{(n-\nu+1)},
\]
auch in den beiden Grenzfällen $\nu=0$, $\nu=(n+1):2$, und somit stimmt also (10) mit dem überein, was sich aus der rechten Seite von (7) durch Vertauschung von $n$ mit $n+1$ ergiebt.

Ganz ebenso ergiebt die Formel (8)
\[
x B_n(x) - B_{n-1}(x) = \overset{\nu}{\Sigma}(-1)^\nu x^{n-2\nu}\left(B_\nu^{(n-\nu-1)} + B_{\nu-1}^{(n-\nu-1)}\right),
\]
und da
\[
B_\nu^{(n-\nu-1)} + B_{\nu-1}^{(n-\nu-1)} = B_\nu^{(n-\nu)}
\]
ist, so ist damit auch (8) allgemein bewiesen.

Der Uebersicht halber setzen wir die Formeln etwas ausführlicher her:
\weq{11}{A_n(x) = x^n - n x^{n-2} + \dfrac{n(n-3)}{1\,.\,2} x^{n-4} - \dfrac{n(n-4)(n-5)}{1\,.\,2\,.\,3} x^{n-6} + \cdots}
\weq{12}{B_n(x) = x^{n-1} - (n-2) x^{n-3} + \dfrac{(n-3)(n-4)}{1\,.\,2} x^{n-5} - \dfrac{(n-4)(n-5)(n-6)}{1\,.\,2\,.\,3} x^{n-7} + \cdots}
Nehmen wir $\cos n\varphi$ und $\sin n\varphi$ als gegeben an, so dient die Gleichung $n^{\text{ten}}$ Grades
\weq{13}{2 \cos n\varphi = A_n(x)}
zur Bestimmung der Unbekannten $x = 2\cos\varphi$. Die Bedeutung der $n$ Wurzeln dieser Gleichung ergiebt sich daraus, dass der Cosinus sich nicht ändert, wenn der Winkel um ein Vielfaches von $2\pi$ wächst. Demnach genügt der Gleichung (13) jeder Werth
\weq{14}{x = 2\cos\left(\varphi + \dfrac{2\nu\pi}{n}\right),}
wenn $\nu$ eine beliebige ganze positive oder negative Zahl bedeutet. Man erhält aber alle von einander verschiedenen Werthe von (14), wenn man $\nu$ ein volles Restsystem in Bezug auf $n$ durchlaufen lässt, also z. B. $\nu = 0, 1, 2 \ldots n-1$ setzt, und man
%CONT>>
